\chapter{Total Protein and Albumin/Globulin Ratio}

\section*{What Is This Test?}
Total protein measures the sum of all serum proteins, of which albumin and globulins together constitute approximately 95\%. The albumin/globulin (A/G) ratio compares albumin to globulin concentrations and aids in classifying disturbances of protein metabolism. A low A/G ratio (with low albumin, high globulin, or both) suggests hepatic synthetic failure, nephrotic syndrome, protein-losing enteropathy, chronic inflammation, or monoclonal gammopathy. A high A/G ratio is rarely clinically significant but may be observed in immune deficiency or low globulin output. Total protein and A/G ratio are part of the comprehensive metabolic panel and a standard liver and immune function workup.

\section*{Alternative Names}
Total Protein; Serum Total Protein; Albumin/Globulin Ratio; A/G Ratio; TP and A/G

\section*{Principle}
Total protein: biuret reaction (cupric ions in alkaline solution react with peptide bonds to form violet color proportional to protein concentration). Albumin: bromocresol green or bromocresol purple dye-binding. Globulin: calculated as total protein minus albumin. A/G ratio calculated as albumin / (total protein $-$ albumin). Performed on automated clinical chemistry analyzers.

\section*{History}
The biuret reaction was described by Riegler in 1914. Dye-binding albumin methods developed in the 1960s with bromocresol green (BCG) and later bromocresol purple (BCP) immunonephelometric methods. The A/G ratio replaced less reliable saline-solubility tests in the mid-20th century.

\section*{Why This Test Is Ordered}
\begin{itemize}
  \item General health screening with CMP
  \item Evaluation of suspected liver disease (cirrhosis, hepatitis)
  \item Investigations of nephrotic syndrome or proteinuria
  \item Suspected multiple myeloma or plasma cell dyscrasias (with elevated globulins)
  \item Chronic inflammatory or autoimmune diseases
  \item Evaluation of edema, ascites, or unexplained weight loss
  \item Monitoring nutritional status in chronic disease
\end{itemize}

\section*{Primary vs Secondary Test}
Primary screening, included in routine CMP alongside albumin.

\section*{How This Test Is Performed}
Venous blood drawn in serum separator or red-top tube; results available within 1--4 hours. Fasting preferred but not strictly required.

\section*{What Preparation Is Needed}
None required; avoid vigorous exercise before blood draw which may transiently raiseglobulins.

\section*{Contraindications and Precautions}
Hemolysis falsely lowers total protein (releasing red cell proteins into plasma). Lipemic or icteric specimens may interfere with photometric measurements. Prolonged tourniquet application causes hemoconcentration.

\section*{Risks}
Venipuncture only.

\section*{What Happens After the Test}
Pattern interpretation: low albumin $+$ low globulin suggests impaired synthesis or loss; low albumin $+$ high globulin suggests cirrhosis or chronic inflammation; normal albumin $+$ high globulin suggests myeloma or polyclonal gammopathy; high albumin $+$ low globulin suggests dehydration or immune deficiency. Abnormal patterns drive further evaluation.

\section*{Understanding Results}

\subsection*{Below Normal}
\begin{itemize}
  \item \textbf{Low total protein (hypoproteinemia):} Malnutrition, liver failure, nephrotic syndrome, protein-losing enteropathy, severe burns, hemorrhage
  \item \textbf{Low albumin (hypoalbuminemia):} Cirrhosis, nephrotic syndrome, protein-losing enteropathy, malnutrition, chronic inflammation, burns, pregnancy
  \item \textbf{Normal/low with high globulin}: Multiple myeloma, Waldenström macroglobulinemia, chronic infection (HIV, hepatitis), autoimmune disease
\end{itemize}

\subsection*{Above Normal}
\begin{itemize}
  \item \textbf{High total protein (hyperproteinemia):} Dehydration (relative), multiple myeloma, polyclonal gammopathy (chronic inflammatory disease)
  \item \textbf{High albumin}: Dehydration (relative); rarely clinically significant elevations
  \item \textbf{High globulin (polyclonal):} Chronic infection, autoimmune disease, liver disease
  \item \textbf{High globulin (monoclonal):} Multiple myeloma, MGUS, Waldenström
\end{itemize}

\section*{Normal Results}
\begin{itemize}
  \item \textbf{Total protein:} 6.0--8.3 g/dL
  \item \textbf{Albumin:} 3.5--5.0 g/dL
  \item \textbf{Globulin (calculated):} 2.0--3.5 g/dL
  \item \textbf{A/G ratio:} 1.1--2.5
  \item \textbf{Low A/G ratio $<$ 1.0:} Impaired hepatic synthesis, nephrotic syndrome, chronic inflammation, multiple myeloma
\end{itemize}

\section*{Follow-up Tests}
\begin{itemize}
  \item \textbf{Serum protein electrophoresis (SPEP):} Identifies polyclonal vs monoclonal gammopathies
  \item \textbf{Immunofixation and serum free light chains:} For suspected myeloma or MGUS
  \item \textbf{Liver function tests (ALT, AST, ALP, bilirubin, PT/INR):} For hepatic involvement
  \item \textbf{Urinalysis and urine protein/creatinine ratio:} For nephrotic syndrome
  \item \textbf{24-hour urine protein and urine protein electrophoresis (UPEP):} Quantify and characterize proteinuria
  \item \textbf{Stool alpha-1 antitrypsin clearance}: Protein-losing enteropathy
  \item \textbf{Fecal calprotectin or imaging}: For inflammatory bowel disease
\end{itemize}

\section*{References}
\begin{enumerate}
  \item Burtis CA, Ashwood ER, Bruns DE. Tietz Textbook of Clinical Chemistry and Molecular Diagnostics. Elsevier; 2023.
  \item Katzmann JA, et al. Screening panels for monoclonal gammopathies. Clin Chem. 2012;58(4):799-803.
  \item Wallach J. Interpretation of Diagnostic Tests. 9th ed. Wolters Kluwer; 2022.
  \item Anderson NL, Anderson NG. The human plasma proteome. Mol Cell Proteomics. 2002;1(11):845-867.
\end{enumerate}

\clearpage
\section*{Regulatory Status and Authorization Date}
Regulatory status applies to the specific assay, instrument, manufacturer, and intended use rather than to the test name alone. No single approval date can be assigned to this test category. For a commercial assay, verify the current product in the relevant regulator's database: the U.S. Food and Drug Administration (FDA) clearance, approval, or authorization; India's Central Drugs Standard Control Organisation (CDSCO) licence or permission; the European Union In Vitro Diagnostic Regulation (EU IVDR) CE certificate issued through the applicable conformity-assessment route; or the United Kingdom Medicines and Healthcare products Regulatory Agency (MHRA) registration and UKCA/CE status. Laboratory-developed tests may not have an individual FDA, CDSCO, EU IVDR, or MHRA approval date.
\clearpage
